\documentclass[11pt,a4paper]{article}
\usepackage[utf8]{inputenc}
\usepackage[T1]{fontenc}
\usepackage[french]{babel}
\usepackage{geometry}
\usepackage{xcolor}
\usepackage{booktabs}
\usepackage{array}
\usepackage{enumitem}
\usepackage{fancyhdr}

\geometry{margin=2cm}
\setlength{\headheight}{14pt}
\pagestyle{fancy}
\fancyhf{}
\rhead{Soutenance Mobile Forensics}
\lfoot{UQAC - Atelier Cyber Sécurité}
\rfoot{\thepage}

\title{\textbf{Répartition des Tâches et Rôles} \\ 
       \large Soutenance Analyse Forensique Samsung Galaxy S10}
\author{Justin Bossard, Mattéo Gouhier, Paul Mathé, Samuel Plet}
\date{Octobre 2025}

\begin{document}

\maketitle

\section{Contexte du Projet}

Dans le cadre du \textbf{DFRWS Challenge 2021}, notre équipe de 4 personnes a réalisé une analyse forensique complète d'un Samsung Galaxy S10. Cette soutenance de 20 minutes présente nos résultats selon un format dynamique et interactif.

\section{Répartition Réelle des Tâches de Recherche}

\subsection{Justin Bossard - \textcolor{blue}{Chef de Projet \& Méthodologie}}
\begin{itemize}[leftmargin=*]
    \item \textbf{Coordination générale} : Mise en forme du rapport final
    \item \textbf{Méthodologie forensique} : Application des 5 phases (Identification → Collection → Examination → Analysis → Presentation)
    \item \textbf{Collaboration avec Samuel} : 
    \begin{itemize}
        \item Identification du suspect (Werner von Strucker)
        \item Analyse de géolocalisation (16 photos en Suisse)
    \end{itemize}
    \item Investigation du portefeuille Bitcoin Cash
    \item \textbf{Intégrité des données} : Vérification SHA2-256
    \item \textbf{Synthèse juridique} : Recevabilité des preuves
\end{itemize}

\subsection{Mattéo Gouhier - \textcolor{orange}{Cryptanalyse \& Messages}}
\begin{itemize}[leftmargin=*]
    \item \textbf{Reverse engineering} : Décompilation APK OmniNotes
    \item \textbf{Cryptanalyse} :
    \begin{itemize}
        \item Identification algorithme AES/ECB + Base64
        \item Découverte préfixe "HHY" (Heil Hydra)
        \item Clé hybride : \texttt{rb5CuefkDLyb9T} + date (YYYY-MM-DD)
    \end{itemize}
    \item \textbf{Déchiffrement messages} :
    \begin{itemize}
        \item 42 messages décodés au total
        \item Messages clés : recrutement (8 avril), rendez-vous (15 avril)
        \item Corrélation "below Turkish consulate" $\leftrightarrow$ photos ICC Genève
    \end{itemize}
    \item \textbf{Développement d'outils} : Scripts de déchiffrement automatique
\end{itemize}

\subsection{Paul Mathé - \textcolor{red}{Audio Forensics \& Skimming}}
\begin{itemize}[leftmargin=*]
    \item \textbf{Analyse audio} :
    \begin{itemize}
        \item Enregistrement suspect : \texttt{2021\_04\_09T1621}
        \item Corrélation temporelle avec date de skimming
        \item Analyse spectrale et patterns audio
    \end{itemize}
    \item \textbf{Investigation skimming} :
    \begin{itemize}
        \item Identification dispositif "audioskim"
        \item Corrélation avec ATM Lausanne (9 avril)
        \item Preuves physiques du crime
    \end{itemize}
    \item \textbf{Expertise technique} : Compréhension mécanismes de skimming
    \item \textbf{Timeline criminelle} : Reconstitution séquence d'événements
\end{itemize}

\subsection{Samuel Plet - \textcolor{green}{OSINT}}
\begin{itemize}[leftmargin=*]
    \item \textbf{Identification du suspect} :
    \begin{itemize}
        \item Corrélation profil Instagram $\leftrightarrow$ photos téléphone
        \item Email suspect : \texttt{therealredskull41@gmail.com}
        \item Lien avec l'univers Marvel/HYDRA
    \end{itemize}
    \item \textbf{Géolocalisation avancée} :
    \begin{itemize}
        \item Analyse des métadonnées EXIF
        \item Cache Google Maps (Bussigny, Prilly, Renens)
        \item Reconstitution du voyage Genève (15 avril)
    \end{itemize}
\end{itemize}

\section{Répartition pour la Soutenance (Format Dynamique)}

\begin{table}[h!]
\centering
\begin{tabular}{|c|p{3cm}|p{5cm}|p{4cm}|}
\hline
\textbf{Acte} & \textbf{Durée} & \textbf{Intervenants} & \textbf{Contenu Principal} \\
\hline
\textbf{Acte 1} & 5 min & \textcolor{blue}{Justin} + \textcolor{green}{Samuel} + \textcolor{red}{Paul} & Contexte + Identité + Lien HYDRA \\
\hline
\textbf{Acte 2} & 7 min & \textcolor{green}{Samuel} + \textcolor{orange}{Mattéo} + \textcolor{red}{Paul} & Géolocalisation + Cryptanalyse + Corrélations \\
\hline
\textbf{Acte 3} & 5 min & \textcolor{blue}{Justin} + \textcolor{red}{Paul} + \textcolor{orange}{Mattéo} & Méthodologie + Audio + Skimming \\
\hline
\textbf{Final} & 3 min & \textcolor{red}{Paul} + \textbf{Tous} & Timeline collaborative + Conclusion \\
\hline
\end{tabular}
\end{table}

\section{Fil Conducteur Dynamique de la Soutenance}

\subsection{ACTE 1 : L'ENQuÊTE COMMENCE (5 min)}
\begin{itemize}[leftmargin=*]
    \item \textbf{\textcolor{blue}{Justin} ouvre (1.5 min)} : Contexte DFRWS, skimming Lausanne → arrestation Genève
    \item \textbf{\textcolor{green}{Samuel} révèle (2 min)} : Identité Werner von Strucker via Autopsy/Instagram
    \item \textbf{\textcolor{red}{Paul} confirme (1.5 min)} : Lien HYDRA, organisation criminelle
    \item \textbf{Transition Samuel} : \textit{"Mattéo, que nous révèlent ses déplacements ?"}
\end{itemize}

\subsection{ACTE 2 : LA TRAQUE (7 min)}
\begin{itemize}[leftmargin=*]
    \item \textbf{\textcolor{orange}{Mattéo} démarre (2.5 min)} : Géolocalisation, cache Google Maps, voyage Genève
    \item \textbf{\textcolor{orange}{Mattéo} révèle (2.5 min)} : Messages chiffrés OmniNotes, préfixe HHY
    \item \textbf{\textcolor{orange}{Mattéo} décode (1.5 min)} : Algorithme AES/ECB, "Heil Hydra", rendez-vous
    \item \textbf{\textcolor{red}{Paul} corrèle (0.5 min)} : "below Turkish consulate" $\leftrightarrow$ photos ICC
    \item \textbf{Transition Paul} : \textit{"Justin, parlons méthodologie !"}
\end{itemize}

\subsection{ACTE 3 : LES PREUVES S'ACCUMULENT (5 min)}
\begin{itemize}[leftmargin=*]
    \item \textbf{\textcolor{blue}{Justin} cadre (1.5 min)} : Principe de Locard, intégrité SHA2-256
    \item \textbf{\textcolor{red}{Paul} dévoile (2 min)} : Bitcoin Cash, réseau d'obfuscation 18 avril
    \item \textbf{\textcolor{red}{Paul} prouve (1.5 min)} : Audio 9 avril, dispositif audioskim ATM
    \item \textbf{Transition Paul} : \textit{"Timeline finale à 4 voix !"}
\end{itemize}

\subsection{FINAL : TIMELINE COLLABORATIVE (3 min)}
\begin{itemize}[leftmargin=*]
    \item \textbf{\textcolor{blue}{Justin}} : \textit{"8 avril : Recrutement HYDRA"}
    \item \textbf{\textcolor{green}{Samuel}} : \textit{"9 avril : Skimming à Lausanne"}
    \item \textbf{\textcolor{orange}{Mattéo}} : \textit{"15 avril : Rendez-vous secret Genève"}
    \item \textbf{\textcolor{red}{Paul}} : \textit{"18 avril : Blanchiment Bitcoin Cash"}
    \item \textbf{Tous ensemble} : \textit{"20 avril : Arrestation - FIN DE PARTIE !"}
    \item \textbf{\textcolor{blue}{Justin} clôture} : Chaîne de preuves, recevabilité juridique
\end{itemize}

\vfill
\centering
\textit{Ce document sert de fil conducteur pour la soutenance dynamique de 20 minutes (prises de parole entremêlées).}

\end{document}